\documentclass[10pt]{extarticle}
\usepackage[a4paper,margin=1.5cm]{geometry}
\usepackage[T1]{fontenc}
\usepackage{amsmath,amssymb}
\usepackage{graphicx}
\usepackage[export]{adjustbox}
\usepackage{multicol}
\usepackage{enumitem}
\usepackage{booktabs,array,tabularx}
\setlength{\columnsep}{0.7cm}
\setlength{\parindent}{0pt}
\setlength{\parskip}{3pt}
\pagestyle{empty}
\makeatletter
% figure registry (figdefs.tex is generated by the extract step)
\newcommand{\qfigdef}[4]{\@namedef{qf@file@#1}{#2}\@namedef{qf@w@#1}{#3}\@namedef{qf@wide@#1}{#4}}
\newcommand{\qf@check}[1]{\@ifundefined{qf@file@#1}{\PackageError{score2}{Unknown figure id #1}{}}{}}
\newcommand{\qf@img}[4][]{\includegraphics[width=#2\dimexpr\@nameuse{qf@w@#4}pt\relax,max width=#3,max height=0.45\textheight,#1]{\@nameuse{qf@file@#4}}}
% Full-width questions: finalize wraps every question that contains a wide figure in \qfullbegin...\qfullend,
% so the WHOLE question (stem, figures, options) is set across the page in reading order, then the two columns
% resume. The question's head - its number, stem and first figure (\qfullsplit marks the end of the head) -
% is kept on one page, so a number is never left at a page foot with its figure overleaf; the rest may break
% across pages normally, so no big gaps are left. Without a \qfullsplit, the whole question is kept together
% if it is under 45% of a page. Never write these yourself.
\newif\ifqf@full
\newif\ifqf@boxed
\newsavebox{\qf@box}
\newcommand{\qf@open}{\global\qf@fulltrue\par\global\qf@boxedtrue
  \setbox\qf@box\vbox\bgroup\hsize\textwidth\linewidth\textwidth}
\newcommand{\qf@emit}[1]{\par\egroup\global\qf@boxedfalse
  \ifdim\dimexpr\ht\qf@box+\dp\qf@box\relax<#1\textheight\noindent\usebox{\qf@box}\par\else\unvbox\qf@box\fi}
\newcommand{\qfullsplit}{\ifqf@boxed\qf@emit{0.85}\fi}
\newcommand{\qf@close}{\ifqf@boxed\qf@emit{0.45}\fi\global\qf@fullfalse}
\newcommand{\qfullbegin}{\end{multicols}\qf@open}
\newcommand{\qfullend}{\qf@close\begin{multicols}{2}}
\newcommand{\qfullstart}{\qf@open}
\newcommand{\qfulldone}{\qf@close}
% \qfig[scale]{id}: block figure, centred, directly where it is written (size set by the tool).
% Fallback if a wide figure is not inside a full-width question: break out of the columns in place - no floats.
\newcommand{\qfig}[2][1]{\qf@check{#2}%
  \ifqf@full\qf@blk{#1}{#2}\else
  \ifnum\@nameuse{qf@wide@#2}=1 \expandafter\@firstoftwo\else\expandafter\@secondoftwo\fi
  {\end{multicols}\begin{center}\qf@img{#1}{\textwidth}{#2}\end{center}\begin{multicols}{2}}%
  {\qf@blk{#1}{#2}}\fi}
\newcommand{\qf@blk}[2]{\par\smallskip{\centering\qf@img{#1}{\linewidth}{#2}\par}\smallskip}
% \qfigi[scale]{id}: inline figure (inside option lists / tables / minipages); never breaks columns.
\newcommand{\qfigi}[2][1]{\qf@check{#2}\qf@img[valign=c]{#1}{\linewidth}{#2}}
\makeatother
% structure macros - lint reads these to check completeness and figure placement
\newcommand{\papersec}[1]{\par\bigskip{\centering\Large\bfseries PAPER-#1\par}\medskip}
\newcommand{\partsec}[2]{\par\medskip\noindent{\large\bfseries PART-#1: #2}\par\smallskip}
\newcommand{\sect}[1]{\par\smallskip\noindent\textbf{#1}\par}
\newcommand{\qn}[1]{\par\smallskip\noindent\textbf{#1.}~}
\newcommand{\paraq}[2]{\par\smallskip\noindent\textbf{Paragraph for Questions #1 and #2}\par}
\newenvironment{qoptions}{\begin{itemize}[leftmargin=1.8em,itemsep=1pt,topsep=2pt,label={}]}{\end{itemize}}
\newcommand{\opt}[2]{\item[(#1)] #2}
% figure registry (generated by score2_tools.py finalize)
\qfigdef{p1_phy_s1_q01_1}{TNAS - Test 06 (30 Jul 2023)-img/p1_phy_s1_q01_1.png}{87.4}{0}
\qfigdef{p1_phy_s1_q02_1}{TNAS - Test 06 (30 Jul 2023)-img/p1_phy_s1_q02_1.png}{106.8}{0}
\qfigdef{p1_phy_s2_q06_1}{TNAS - Test 06 (30 Jul 2023)-img/p1_phy_s2_q06_1.png}{236.6}{0}
\qfigdef{p1_phy_s2_q07_1}{TNAS - Test 06 (30 Jul 2023)-img/p1_phy_s2_q07_1.png}{84.5}{0}
\qfigdef{p1_phy_s3_q08_1}{TNAS - Test 06 (30 Jul 2023)-img/p1_phy_s3_q08_1.png}{100.2}{0}
\qfigdef{p1_phy_s3_q09_1}{TNAS - Test 06 (30 Jul 2023)-img/p1_phy_s3_q09_1.png}{108.3}{0}
\qfigdef{p1_phy_s3_q10_1}{TNAS - Test 06 (30 Jul 2023)-img/p1_phy_s3_q10_1.png}{167.3}{0}
\qfigdef{p1_phy_s3_q11_1}{TNAS - Test 06 (30 Jul 2023)-img/p1_phy_s3_q11_1.png}{88.9}{0}
\qfigdef{p1_phy_s3_q11_2}{TNAS - Test 06 (30 Jul 2023)-img/p1_phy_s3_q11_2.png}{112.7}{0}
\qfigdef{p1_phy_s3_q11_3}{TNAS - Test 06 (30 Jul 2023)-img/p1_phy_s3_q11_3.png}{117.3}{0}
\qfigdef{p1_phy_s3_q11_4}{TNAS - Test 06 (30 Jul 2023)-img/p1_phy_s3_q11_4.png}{101.4}{0}
\qfigdef{p1_phy_s4_q02_1}{TNAS - Test 06 (30 Jul 2023)-img/p1_phy_s4_q02_1.png}{94.2}{0}
\qfigdef{p1_phy_s4_q03_1}{TNAS - Test 06 (30 Jul 2023)-img/p1_phy_s4_q03_1.png}{150.3}{0}
\qfigdef{p1_phy_s4_q04_1}{TNAS - Test 06 (30 Jul 2023)-img/p1_phy_s4_q04_1.png}{102.7}{0}
\qfigdef{p1_phy_s4_q05_1}{TNAS - Test 06 (30 Jul 2023)-img/p1_phy_s4_q05_1.png}{54.6}{0}
\qfigdef{p1_phy_s4_q06_1}{TNAS - Test 06 (30 Jul 2023)-img/p1_phy_s4_q06_1.png}{145.4}{0}
\qfigdef{p1_chem_s1_q03_1}{TNAS - Test 06 (30 Jul 2023)-img/p1_chem_s1_q03_1.png}{32.0}{0}
\qfigdef{p1_chem_s1_q04_1}{TNAS - Test 06 (30 Jul 2023)-img/p1_chem_s1_q04_1.png}{323.5}{1}
\qfigdef{p1_chem_s1_q04_A}{TNAS - Test 06 (30 Jul 2023)-img/p1_chem_s1_q04_A.png}{64.6}{0}
\qfigdef{p1_chem_s1_q04_B}{TNAS - Test 06 (30 Jul 2023)-img/p1_chem_s1_q04_B.png}{71.3}{0}
\qfigdef{p1_chem_s1_q04_C}{TNAS - Test 06 (30 Jul 2023)-img/p1_chem_s1_q04_C.png}{73.9}{0}
\qfigdef{p1_chem_s2_q07_1}{TNAS - Test 06 (30 Jul 2023)-img/p1_chem_s2_q07_1.png}{92.5}{0}
\qfigdef{p1_chem_s3_q10_1}{TNAS - Test 06 (30 Jul 2023)-img/p1_chem_s3_q10_1.png}{107.8}{0}
\qfigdef{p1_chem_s3_q11_1}{TNAS - Test 06 (30 Jul 2023)-img/p1_chem_s3_q11_1.png}{34.3}{0}
\qfigdef{p1_chem_s3_q11_2}{TNAS - Test 06 (30 Jul 2023)-img/p1_chem_s3_q11_2.png}{34.2}{0}
\qfigdef{p1_chem_s3_q11_3}{TNAS - Test 06 (30 Jul 2023)-img/p1_chem_s3_q11_3.png}{18.6}{0}
\qfigdef{p1_chem_s3_q11_4}{TNAS - Test 06 (30 Jul 2023)-img/p1_chem_s3_q11_4.png}{37.0}{0}
\qfigdef{p1_chem_s4_q04_1}{TNAS - Test 06 (30 Jul 2023)-img/p1_chem_s4_q04_1.png}{58.5}{0}
\qfigdef{p1_chem_s4_q04_2}{TNAS - Test 06 (30 Jul 2023)-img/p1_chem_s4_q04_2.png}{50.2}{0}
\qfigdef{p2_phy_s1_q01_1}{TNAS - Test 06 (30 Jul 2023)-img/p2_phy_s1_q01_1.png}{53.1}{0}
\qfigdef{p2_phy_s1_q02_1}{TNAS - Test 06 (30 Jul 2023)-img/p2_phy_s1_q02_1.png}{111.1}{0}
\qfigdef{p2_phy_s1_q02_2}{TNAS - Test 06 (30 Jul 2023)-img/p2_phy_s1_q02_2.png}{129.4}{0}
\qfigdef{p2_phy_s1_q04_1}{TNAS - Test 06 (30 Jul 2023)-img/p2_phy_s1_q04_1.png}{118.8}{0}
\qfigdef{p2_phy_s1_q04_A}{TNAS - Test 06 (30 Jul 2023)-img/p2_phy_s1_q04_A.png}{71.3}{0}
\qfigdef{p2_phy_s1_q04_B}{TNAS - Test 06 (30 Jul 2023)-img/p2_phy_s1_q04_B.png}{71.5}{0}
\qfigdef{p2_phy_s1_q04_C}{TNAS - Test 06 (30 Jul 2023)-img/p2_phy_s1_q04_C.png}{79.3}{0}
\qfigdef{p2_phy_s1_q04_D}{TNAS - Test 06 (30 Jul 2023)-img/p2_phy_s1_q04_D.png}{74.2}{0}
\qfigdef{p2_phy_s2_q05_1}{TNAS - Test 06 (30 Jul 2023)-img/p2_phy_s2_q05_1.png}{108.3}{0}
\qfigdef{p2_phy_s2_q06_1}{TNAS - Test 06 (30 Jul 2023)-img/p2_phy_s2_q06_1.png}{76.4}{0}
\qfigdef{p2_phy_s2_q07_1}{TNAS - Test 06 (30 Jul 2023)-img/p2_phy_s2_q07_1.png}{142.7}{0}
\qfigdef{p2_phy_s4_q05_1}{TNAS - Test 06 (30 Jul 2023)-img/p2_phy_s4_q05_1.png}{87.4}{0}
\qfigdef{p2_phy_s4_q06_1}{TNAS - Test 06 (30 Jul 2023)-img/p2_phy_s4_q06_1.png}{82.0}{0}
\qfigdef{p2_phy_s4_q07_1}{TNAS - Test 06 (30 Jul 2023)-img/p2_phy_s4_q07_1.png}{92.6}{0}
\qfigdef{p2_phy_s4_q08_1}{TNAS - Test 06 (30 Jul 2023)-img/p2_phy_s4_q08_1.png}{107.0}{0}
\qfigdef{p2_phy_s4_q09_1}{TNAS - Test 06 (30 Jul 2023)-img/p2_phy_s4_q09_1.png}{160.9}{0}
\qfigdef{p2_phy_s4_q10_1}{TNAS - Test 06 (30 Jul 2023)-img/p2_phy_s4_q10_1.png}{100.9}{0}
\qfigdef{p2_chem_s1_q04_1}{TNAS - Test 06 (30 Jul 2023)-img/p2_chem_s1_q04_1.png}{97.9}{0}
\qfigdef{p2_chem_s2_q06_1}{TNAS - Test 06 (30 Jul 2023)-img/p2_chem_s2_q06_1.png}{266.2}{1}
\qfigdef{p2_chem_s2_q06_2}{TNAS - Test 06 (30 Jul 2023)-img/p2_chem_s2_q06_2.png}{55.6}{0}
\qfigdef{p2_chem_s2_q07_A}{TNAS - Test 06 (30 Jul 2023)-img/p2_chem_s2_q07_A.png}{97.0}{0}
\qfigdef{p2_chem_s2_q07_B}{TNAS - Test 06 (30 Jul 2023)-img/p2_chem_s2_q07_B.png}{122.2}{0}
\qfigdef{p2_chem_s2_q07_C}{TNAS - Test 06 (30 Jul 2023)-img/p2_chem_s2_q07_C.png}{105.6}{0}
\qfigdef{p2_chem_s3_para1_2_1}{TNAS - Test 06 (30 Jul 2023)-img/p2_chem_s3_para1_2_1.png}{135.9}{0}
\qfigdef{p2_chem_s4_q09_1}{TNAS - Test 06 (30 Jul 2023)-img/p2_chem_s4_q09_1.png}{108.5}{0}
\qfigdef{p2_chem_s4_q09_2}{TNAS - Test 06 (30 Jul 2023)-img/p2_chem_s4_q09_2.png}{245.2}{0}
\raggedcolumns
\makeatletter\@beginparpenalty=10000\makeatother
\newcolumntype{L}[1]{>{\raggedright\arraybackslash}p{#1}}
\begin{document}
\begin{center}{\Large\bfseries TNAS --- Test 06 (30 Jul 2023)}\end{center}
\begin{multicols}{2}
\raggedcolumns
\papersec{1}
\partsec{1}{PHYSICS}
\sect{SECTION-I(i)}
\qn{1} Water in a clean aquarium forms a meniscus, as illustrated in the figure. The difference in height $h$ between the centre and the edge of the meniscus will be (Surface tension of water $=0.02$ N/m, contact angle $=0^\circ$, density of water $=1000$ kg/m$^3$, $g=10$ m/s$^2$)
\qfig{p1_phy_s1_q01_1}
\begin{qoptions}
\opt{A}{1 mm}
\opt{B}{2 mm}
\opt{C}{3 mm}
\opt{D}{4 mm}
\end{qoptions}
\qn{2} Two straight and very long non-magnetic conductors $A$ and $B$, insulated from each other, carry equal and opposite current $I$ in the positive and the negative $z$-direction, respectively. The cross-sections of the conductors (shaded in the figure) are in the $x$-$y$ plane as shown in figure with a distance $R$ between their centres. The current $I$ in each conductor is uniformly distributed over the shaded cross-section. The resulting cross-section of each have an area $\dfrac{\pi R^2}{3}+\dfrac{\sqrt3R^2}{2}$. The magnetic field $B$ at any point on $y$-axis in the space between the conductors is :-
\qfig{p1_phy_s1_q02_1}
\begin{qoptions}
\opt{A}{$\dfrac{9\mu_0I}{\left(2\pi+3\sqrt3\right)R}$}
\opt{B}{$\dfrac{3\mu_0I}{2\left(2\pi+3\sqrt3\right)R}$}
\opt{C}{$\dfrac{3\mu_0I}{\left(2\pi+3\sqrt3\right)R}$}
\opt{D}{$\dfrac{6\mu_0I}{\left(2\pi+3\sqrt3\right)R}$}
\end{qoptions}
\qn{3} Choose the CORRECT option for electric field components of unpolarized light passing through polarizing sheet.
\begin{qoptions}
\opt{A}{An electric field component parallel to polarizing direction is absorbed and perpendicular to polarizing direction get transmitted.}
\opt{B}{An electric field component parallel to polarizing direction is transmitted and perpendicular to polarizing direction get absorbed.}
\opt{C}{An electric field component parallel to polarizing direction is rotated through some angle and perpendicular component is transmitted.}
\opt{D}{An electric field component parallel to polarizing direction is transmitted and perpendicular component get rotated through some angle.}
\end{qoptions}
\qn{4} Consider a set of two slits of width $b=5\times10^{-2}$ cm and separated by a distance of $d=0.1$ cm. They are illuminated by light of wavelength 600 nm. If a convex lens of focal length 10 cm is placed beyond the slits and resulting interference pattern is observed on a screen in the focal plane of lens :-
\begin{qoptions}
\opt{A}{Within the 1st diffraction minima, we observe 3 maxima due to interference.}
\opt{B}{Within the 1st diffraction minima, we observe 5 maxima due to interference.}
\opt{C}{A maxima within the 1st diffraction minima is located at $y=0.03$ mm.}
\opt{D}{A maxima within the 1st diffraction minima is located at $y=0.09$ mm.}
\end{qoptions}
\sect{SECTION-I(ii)}
\qn{5} A standing wave is setup in a string fixed at both ends. Then
\begin{qoptions}
\opt{A}{the sum of total energy per unit length at nodes and antinodes is constant}
\opt{B}{the total energy of a point midway between a node and an adjacent antinode is constant}
\opt{C}{if the string has a standing wave and a component of traveling wave, then the kinetic energy of the points of minimum amplitude is non-zero at some instant of time.}
\opt{D}{if the string has a standing wave and a component of traveling wave, then the potential energy of the points of maximum amplitude is non zero at some instant of time.}
\end{qoptions}
\qn{6} A beam of light traveling horizontally consists of an unpolarized component with intensity $I_0$ and a polarized component with intensity $I_p$. The plane of polarization is oriented at an angle $\theta$ with respect to the vertical. The figure shows the total intensity $I_{\text{total}}$ after the light passes through a polarizer as a function of the angle $\alpha$, that the pass axis of the polarizer makes with respect to the vertical. Identify the correct statement (s).
\qfig{p1_phy_s2_q06_1}
\begin{qoptions}
\opt{A}{$\theta=125^\circ$}
\opt{B}{$I_p=20$ W/m$^2$}
\opt{C}{$I_0=17.5$ W/m$^2$}
\opt{D}{$I_0=10$ W/m$^2$}
\end{qoptions}
\qn{7} Two identical polaroid films are pasted on both the slits of YDSE apparatus. $O$ is the centre of the screen. Unpolarized light is incident on the apparatus. The intensity at $O$ is seen to be $I_0$, when no films are pasted. Choose the correct statement(s).
\qfig{p1_phy_s2_q07_1}
\begin{qoptions}
\opt{A}{If pass axes of both films are parallel to each other, intensity at $O$ will be $\dfrac{I_0}{2}$.}
\opt{B}{If pass axes are perpendicular to each other, intensity at $O$ is $\dfrac{I_0}{4}$.}
\opt{C}{If pass axes make an angle of $60^\circ$ with each other, intensity at $O$ is $\dfrac{I_0(\sqrt3+1)}{8}$.}
\opt{D}{If pass axes make an angle of $60^\circ$ with each other, intensity at $O$ is $\dfrac{3I_0}{8}$.}
\end{qoptions}
\sect{SECTION-I(iii)}
\qn{8} A conducting loop moves in $xy$-plane with velocity $v\hat i$ as shown. A uniform magnetic field $\vec B_1$ exists in whole region and a uniform magnetic field $\vec B_2$ exists in addition for all $x>0$. (At shown instant)
\qfig{p1_phy_s3_q08_1}
{\footnotesize\setlength{\tabcolsep}{2pt}\begin{tabular}{|c|p{0.38\linewidth}|c|p{0.3\linewidth}|}
\hline
\multicolumn{2}{|c|}{\textbf{List-I}} & \multicolumn{2}{c|}{\textbf{List-II}}\\ \hline
(P) & $\vec B_1=2\hat i+2\hat j+4\hat k$, $\vec B_2=0$ & (1) & potential $V_Q>V_P$\\ \hline
(Q) & $\vec B_1=\hat i+3\hat j-2\hat k$, $\vec B_2=0$ & (2) & potential $V_P>V_Q$\\ \hline
(R) & $\vec B_1=0$, $\vec B_2=4\hat i+2\hat j+9\hat k$ & (3) & no induced current\\ \hline
(S) & $\vec B_1=0$, $\vec B_2=-\hat i-2\hat j-2\hat k$ & (4) & induced current flows from $Q$ to $P$\\ \hline
 & & (5) & induced current flows from $P$ to $Q$\\ \hline
\end{tabular}}\par
\begin{qoptions}
\opt{A}{$P\to1,3;Q\to2,3;R\to1,2;S\to1,4$}
\opt{B}{$P\to2,3;Q\to1,3;R\to2,5;S\to1,4$}
\opt{C}{$P\to2,3;Q\to2,5;R\to1,4;S\to1,2$}
\opt{D}{$P\to1,2;Q\to4,5;R\to2,3;S\to5$}
\end{qoptions}
\qn{9} A Young's double slit apparatus is shown in figure :-
\qfig{p1_phy_s3_q09_1}
{\footnotesize\setlength{\tabcolsep}{2pt}\begin{tabular}{|c|p{0.38\linewidth}|c|p{0.3\linewidth}|}
\hline
\multicolumn{2}{|c|}{\textbf{List-I}} & \multicolumn{2}{c|}{\textbf{List-II}}\\ \hline
(P) & Source $S$ move upwards & (1) & Central fringe moves upwards\\ \hline
(Q) & A thin glass slab is placed in front of slit $S_2$ & (2) & Central fringe moves downward\\ \hline
(R) & Whole apparatus is immersed in water & (3) & Fringe width remains constant\\ \hline
(S) & Diameter of slit $S_2$ is doubled & (4) & Fringe width decreases\\ \hline
 & & (5) & Intensity of central maxima increases\\ \hline
\end{tabular}}\par
\begin{qoptions}
\opt{A}{$P\to2,3;Q\to2,3;R\to4;S\to3,5$}
\opt{B}{$P\to2,3;Q\to2;R\to4;S\to5$}
\opt{C}{$P\to2,3;Q\to3,5;R\to3,4;S\to5$}
\opt{D}{$P\to3;Q\to2,3;R\to5;S\to3,5$}
\end{qoptions}
\qn{10} Two point monochromatic sources of light are free to move in a circle of radius 1 mm in anticlockwise direction without any external torque with angular velocity of $\pi/4$ rad/sec. In list I various times are given and in list II the coordinate is given (here zeroth maxima / minima is a point where path difference is zero and last maxima/minima is the point when path difference is maximum). Given wavelength of light used $=5000$ \AA\ and neglect Doppler's effect
\qfig{p1_phy_s3_q10_1}
{\footnotesize\setlength{\tabcolsep}{2pt}\begin{tabular}{|c|p{0.35\linewidth}|c|p{0.3\linewidth}|}
\hline
\multicolumn{2}{|c|}{\textbf{List-I}} & \multicolumn{2}{c|}{\textbf{List-II}}\\ \hline
(P) & $t=0$, zeroth maxima is at & (1) & (1 m, 0)\\ \hline
(Q) & $t=0$, last maxima is at & (2) & (1 m, $-\infty$)\\ \hline
(R) & $t=2$, first maxima is at & (3) & (1 m, 0.25 m)\\ \hline
(S) & $t=2$, thousandth maxima is at & (4) & (1 m, 250 $\mu$m)\\ \hline
 & & (5) & None of these\\ \hline
\end{tabular}}\par
\begin{qoptions}
\opt{A}{$P\to1;Q\to2;R\to3;S\to5$}
\opt{B}{$P\to1;Q\to2;R\to3;S\to4$}
\opt{C}{$P\to2;Q\to1;R\to4;S\to5$}
\opt{D}{$P\to2;Q\to3;R\to1;S\to5$}
\end{qoptions}
\qn{11} List-I contains some circuits and List-II contains some properties. In List-II current signifies instantaneous current and potential difference refers instantaneous potential difference.
{\footnotesize\setlength{\tabcolsep}{2pt}\begin{tabular}{|c|p{0.46\linewidth}|c|p{0.25\linewidth}|}
\hline
\multicolumn{2}{|c|}{\textbf{List-I}} & \multicolumn{2}{c|}{\textbf{List-II}}\\ \hline
(P) & \qfigi{p1_phy_s3_q11_1}\par A closely separated conducting wire is wrapped over an insulating cylinder. Wire is connected to a resistance $R$ and an alternating source via a switch $S$, which is closed at $t=0$ & (1) & In branch $x$, current varies\\ \hline
(Q) & \qfigi{p1_phy_s3_q11_2}\par A charged metallic sphere of radius $r$ is earthed with a resistance $R$ and a switch $S$, which is closed at $t=0$. & (2) & Between points $a$ and $b$ magnitude of potential difference increases always.\\ \hline
(R) & \qfigi{p1_phy_s3_q11_3}\par An inductor $L$ and resistor $R$ is connected to a cell of emf $\varepsilon$ and internal resistance $r$ with a switch($s$), which is closed at $t=0$ & (3) & In branch $x$, current is constant\\ \hline
(S) & \qfigi{p1_phy_s3_q11_4}\par A circuit (used to measure an unknown resistance) is in balanced conditions (Galvanometer reads zero) & (4) & Between points $a$ and $b$ potential difference is zero always.\\ \hline
 & & (5) & Rate of heat generation through the resistance $R$ is constant\\ \hline
\end{tabular}}\par
\begin{qoptions}
\opt{A}{$P\to1;Q\to1;R\to1;S\to3,5$}
\opt{B}{$P\to2;Q\to4;R\to3;S\to1$}
\opt{C}{$P\to1;Q\to2;R\to3,4;S\to2$}
\opt{D}{$P\to3;Q\to4;R\to1,2;S\to3,4$}
\end{qoptions}
\sect{SECTION-II}
\qn{1} One end of a string of length $L=15$m is tied to the ceiling of an elevator accelerating upwards with an acceleration 20 m/s$^2$. The other end of the string is free. The linear mass density of the string varies linearly from 0 to $\lambda_0$ from bottom to top. Find the time taken by a pulse (in seconds) to reach from bottom to the loop. ($g=10$ ms$^{-2}$)
\qn{2} A non-conducting ring of radius $R$ having uniformly distributed charge $Q$ starts rotating about $xx'$ axis passing through the centre in the plane of the ring with an angular acceleration $\alpha$ as shown in the figure. Another small conducting ring having radius $a(a\ll R)$ is kept fixed at the centre of bigger ring in such away that the axis $xx'$ is passing through its centre and perpendicular to its plane. If the resistance of the small ring is $r=1\,\Omega$, find the induced current in it (in ampere)\par
(Given $q=\dfrac{16\times10^2}{\mu_0}$C, $R=1$m, $a=0.1$m, $\alpha=8$ rad/s$^2$)
\qfig{p1_phy_s4_q02_1}
\qn{3} The Fresnel biprism is a variation in the YDSE. The fresnel biprism consists of two thin prism joint at their bases. Consider a Fresnal biprism experiment (shown in figure) in which prism angle $2\times10^{-3}$ rad, distance of source (of wavelength 600 nm) from prism is 1m and distance between screen and prism is 4m. Find ratio of maximum intensity on screen to intensity just below point $A$ in interference region. [Assuming intensity of light from upper and lower prism is same in interference region and Refractive index of prism $=1.5$]
\qfig{p1_phy_s4_q03_1}
\qn{4} Refer to the shown diagram. A plane mirror of length $D$, a point source $S$ and a screen are placed as shown. The length $PQ$ on the screen where interference pattern is observed is $\alpha$ mm. The value of $2\alpha$ is
\qfig{p1_phy_s4_q04_1}
\qn{5} A thin film of plastic ($n=1.56$) is 0.25 $\mu$m thick. It is sandwiched between two glass slabs with refractive indices of 1.58 and 1.52, respectively. White light (400 -- 700 nm) is first incident normally on the slab for which $n=1.58$. If $\lambda$ is the visible wavelength (in nm) which is missing in the reflected light, then find the value of $\dfrac\lambda{260}$ in nm.
\qfig{p1_phy_s4_q05_1}
\qn{6} Find the damping constant $c$ (in kg/sec) that gives critical damping of the rigid bar executing motions near the equilibrium position shown in Figure. (Given mass $=6$ kg, $k=200$ N/m, $L=2a$)
\qfig{p1_phy_s4_q06_1}
\partsec{2}{CHEMISTRY}
\sect{SECTION-I(i)}
\qn{1} Which of the following reagent is used with KI solution to give pure HI?
\begin{qoptions}
\opt{A}{conc. $\mathrm{H_2SO_4}$}
\opt{B}{dilute $\mathrm{H_2SO_4}$}
\opt{C}{conc. $\mathrm{H_3PO_4}$}
\opt{D}{dilute HCl}
\end{qoptions}
\qn{2} In a hydrogen like species , an electron jumps from an orbital having two radial and two angular nodes to the orbital having same sign of wave function in all direction, at any distance. If the energy of emitted photon is 326.4 eV, the atom or ion is
\begin{qoptions}
\opt{A}{H}
\opt{B}{$\mathrm{Li^{2+}}$}
\opt{C}{$\mathrm{He^+}$}
\opt{D}{$\mathrm{B^{4+}}$}
\end{qoptions}
\qn{3} \qfigi{p1_chem_s1_q03_1}$\xrightarrow{\text{N.B.S.}}(P)\xrightarrow[\text{ether}]{\text{Mg}}(Q)\xrightarrow[\mathrm{H_2O/H^+}]{\mathrm{CO_2}}(R)\xrightarrow{\text{L.A.H.}}(S)\xrightarrow{\text{PCC}}(T)$\par
Which of the following is correct?
\begin{qoptions}
\opt{A}{D.U. of $(R)$ is 2}
\opt{B}{$(S)$ is $2^\circ$ alcohol}
\opt{C}{$(T)$ have four $sp^3$ hybrid carbon.}
\opt{D}{Reaction of $(R)$ with sodalime gives methyl cyclohexene.}
\end{qoptions}
\qfullbegin
\qn{4}
\qfig{p1_chem_s1_q04_1}
\qfullsplit
$\mathrm{PhCOOH}+\mathrm{CHI_3}\downarrow$ (yellow ppt)\par
The structure of compound $P$ is :
\begin{qoptions}
\opt{A}{\qfigi{p1_chem_s1_q04_A}}
\opt{B}{\qfigi{p1_chem_s1_q04_B}}
\opt{C}{\qfigi{p1_chem_s1_q04_C}}
\opt{D}{$\mathrm{Ph-C\equiv C-Ph}$}
\end{qoptions}
\qfullend
\sect{SECTION-I(ii)}
\qn{5} The vapour pressure of solid benzoic acid has been found to obey the relationship (in the neighbourhood of 298 K) as : $\ln\dfrac{P}{P^0}=a-\dfrac bT$ ; where $a=20$ and $b=1\times10^4$ K. Which of the following values are correct at 298 K for the sublimation of benzoic acid : ($R=8.3$ J/mole-K)
\begin{qoptions}
\opt{A}{$\Delta G^\circ=33.53$ kJ mol$^{-1}$}
\opt{B}{$\Delta H^\circ=83$ kJ mol$^{-1}$}
\opt{C}{$\Delta S^\circ=166$ JK$^{-1}$ mol$^{-1}$}
\opt{D}{$\Delta H^\circ=-88.96$ kJ mol$^{-1}$}
\end{qoptions}
\qn{6} Which is/are correct statement -
\begin{qoptions}
\opt{A}{$\mathrm{H_2(g)}+\frac12\mathrm{O_2(g)}\to\mathrm{H_2O(l)}$ ; $\Delta_rH$\par $\Delta_rH=\Delta H_C(\mathrm{H_2,g})=\Delta H_f(\mathrm{H_2O},\ell)$}
\opt{B}{$\mathrm{Na(s)}\to\mathrm{Na(g)}$ ; $\Delta_rH$\par $\Delta_rH=\Delta H_{\text{sublimation}}(\mathrm{Na,s})=\Delta H_{\text{atomisation}}(\mathrm{Na,s})$}
\opt{C}{$\mathrm{C(graphite)}+\mathrm{O_2(g)}\to\mathrm{CO_2(g)}$ ; $\Delta_rH$\par $\Delta_rH=\Delta H_C(\mathrm{C,graphite})=\Delta H_f(\mathrm{CO_2,g})$}
\opt{D}{$\mathrm{H_2(g)}\to2\mathrm{H(g)}$ ; $\Delta_rH$\par $\Delta_rH=\Delta H_{\text{atomisation}}(\mathrm{H_2,g})=\Delta H_f(\mathrm{H,g})$}
\end{qoptions}
\qn{7} 1 moles of an ideal gas, [$\gamma=1.5$] are taken through a series of processes.
\qfig{p1_chem_s2_q07_1}
Information -I : The temperature at state `A' is 300 K.\par
Information - II : The reversible isothermal expansion from $B$ to $C$ doubles the volume\par
Information - III : The entropy change of the system from $C$ to $D$ is $4\ell n(16)$ Cal/ K.\par
Select the correct option (s) using above information :\par
[$R=2$ cal/mol-K]
\begin{qoptions}
\opt{A}{The entropy change of the system from $D$ to $A$ is $-3R\ \ell n\,8$}
\opt{B}{The temperature at point $D$ is 2400 K}
\opt{C}{The work done from $A$ to $B$ is $-600$ Cal.}
\opt{D}{The work done from $B$ to $C$ is $-300\ \ell n\,2$ Cal.}
\end{qoptions}
\sect{SECTION-I(iii)}
\qn{8} Match the List-I with List-II.\par
{\footnotesize\setlength{\tabcolsep}{2pt}\begin{tabular}{|c|p{0.38\linewidth}|c|p{0.3\linewidth}|}
\hline
\multicolumn{2}{|c|}{\textbf{List-I}} & \multicolumn{2}{c|}{\textbf{List-II}}\\ \hline
(P) & $\mathrm{AlF_3}$, $\mathrm{SF_4}$, $\mathrm{PF_5}$ & (1) & $1^{\text{st}}$ compound has highest boiling point\\ \hline
(Q) & $(\mathrm{CH_3})_2\mathrm{O}$, $\mathrm{S(CH_3)_2}$, $\mathrm{Se(CH_3)_2}$ & (2) & III$^{\text{rd}}$ compound has highest boiling point\\ \hline
(R) & HCl, HF, HBr & (3) & II$^{\text{nd}}$ compound has highest boiling point\\ \hline
(S) & $\mathrm{CH_4}$, $\mathrm{CH_3Cl}$, $\mathrm{CH_3OH}$, $\mathrm{CH_3CO_2H}$ & (4) & IV$^{\text{th}}$ compound has highest boiling point\\ \hline
 & & (5) & III$^{\text{rd}}$ compound has very-very low boiling point\\ \hline
\end{tabular}}\par
The correct option is :
\begin{qoptions}
\opt{A}{$P\to1;Q\to2;R\to3;S\to4$}
\opt{B}{$P\to1;Q\to2;R\to4;S\to3$}
\opt{C}{$P\to2;Q\to1;R\to3;S\to4$}
\opt{D}{$P\to2;Q\to1;R\to4;S\to3$}
\end{qoptions}
\qn{9} Match the List-I with List-II.\par
{\footnotesize\setlength{\tabcolsep}{2pt}\begin{tabular}{|c|p{0.3\linewidth}|c|p{0.38\linewidth}|}
\hline
\multicolumn{2}{|c|}{\textbf{List-I}} & \multicolumn{2}{c|}{\textbf{List-II}}\\ \hline
(P) & Down's cell & (1) & Fused $\mathrm{MgCl_2}$\\ \hline
(Q) & Dow sea water process & (2) & Fused $\mathrm{Al_2O_3}+\mathrm{Na_3AlF_6}+\mathrm{CaF_2}$\\ \hline
(R) & Hall - Heroult & (3) & Fused $\mathrm{KHF_2}$\\ \hline
(S) & Moissan process & (4) & Fused NaCl and $\mathrm{CaCl_2}$\\ \hline
 & & (5) & Fused $\mathrm{BaCl_2}$ and NaCl\\ \hline
\end{tabular}}\par
The correct option is :
\begin{qoptions}
\opt{A}{$P\to4;Q\to3;R\to2;S\to1$}
\opt{B}{$P\to4;Q\to1;R\to3;S\to2$}
\opt{C}{$P\to4;Q\to1;R\to2;S\to3$}
\opt{D}{$P\to4;Q\to3;R\to1;S\to2$}
\end{qoptions}
\qn{10} One mole of a monatomic ideal gas is taken along two cyclic processes $E\to F\to G\to E$ and $E\to F\to H\to E$ as shown in the $p$-$V$ diagram.
\qfig{p1_chem_s3_q10_1}
The processes involved are purely isochoric, isobaric, isothermal or adiabatic. Match the paths in List-I with magnitudes of the work done in List-II and select the correct answer using the codes given below the lists.\par
{\footnotesize\setlength{\tabcolsep}{2pt}\begin{tabular}{|c|p{0.2\linewidth}|c|p{0.3\linewidth}|}
\hline
\multicolumn{2}{|c|}{\textbf{List-I}} & \multicolumn{2}{c|}{\textbf{List-II}}\\ \hline
(P) & $G\to E$ & (1) & $160\,p_0V_0\ln2$\\ \hline
(Q) & $G\to H$ & (2) & $36\,p_0V_0$\\ \hline
(R) & $F\to H$ & (3) & $24\,p_0V_0$\\ \hline
(S) & $F\to G$ & (4) & $31\,p_0V_0$\\ \hline
 & & (5) & $96\,p_0V_0$\\ \hline
\end{tabular}}\par
\begin{qoptions}
\opt{A}{$P\to4;Q\to3;R\to5;S\to2$}
\opt{B}{$P\to3;Q\to2;R\to4;S\to1$}
\opt{C}{$P\to4;Q\to3;R\to2;S\to1$}
\opt{D}{$P\to1;Q\to2;R\to3,5;S\to4$}
\end{qoptions}
\qn{11} Match the list\par
{\footnotesize\setlength{\tabcolsep}{2pt}\begin{tabular}{|c|p{0.2\linewidth}|c|p{0.5\linewidth}|}
\hline
\multicolumn{2}{|c|}{\textbf{List-I}} & \multicolumn{2}{c|}{\textbf{List-II}}\\ \hline
(P) & \qfigi{p1_chem_s3_q11_1} & (1) & Gives orange colour with Brady's reagent (2,4-DNP)\\ \hline
(Q) & \qfigi{p1_chem_s3_q11_2} & (2) & Gives Haloform reaction\\ \hline
(R) & \qfigi{p1_chem_s3_q11_3} & (3) & Can give Cannizaro reaction with conc. NaOH.\\ \hline
(S) & \qfigi{p1_chem_s3_q11_4} & (4) & Can give aldol reaction with dil. NaOH\\ \hline
 & & (5) & Gives white ppt. with alcoholic $\mathrm{AgNO_3}$\\ \hline
\end{tabular}}\par
\begin{qoptions}
\opt{A}{$P\to1,2,4;Q\to2,5;R\to1,4;S\to5$}
\opt{B}{$P\to2,4;Q\to1,3;R\to2,5;S\to4,5$}
\opt{C}{$P\to1,5;Q\to3,4;R\to2,3;S\to1,2$}
\opt{D}{$P\to2,3,4;Q\to4,5;R\to1,2,3;S\to1,4,5$}
\end{qoptions}
\sect{SECTION-II}
\qn{1} How many elements can be purified by distillation process?\par
Zn , Pb , Cd , Hg , Mg , Na , K
\qn{2} Find the number of reagent(s) which oxidises $\mathrm{Br^-}$ ion solution.\par
$\mathrm{MnO_4^{\ominus}/H^{\oplus}}$ , $\mathrm{Cr_2O_7^{2-}/H^{\oplus}}$ , $\mathrm{Cl_2}$--water , $\mathrm{Br_2}$--water , $\mathrm{I_2}$--water
\qn{3} An element $X$ (Atomic weight $=24$ gm/mole) form FCC lattice. If edge length of unit cell is $4\times10^{-8}$ cm and observed density is $2.4\times10^3$ kg/m$^3$, then percentage occupancy of lattice points of $X$ element is (Given $N_A=6\times10^{23}$)
\qn{4} (a) Find out the total compounds of molecular formula $\mathrm{C_6H_{10}}$ which react with 1 eq. of $\mathrm{H_2}$/Pt to produce 1,2-dimethyl cyclobutane, without counting stereoisomers.\par
(b) Total number of structural isomers of $\mathrm{C_4H_{10}O}$.\par
\qfigi{p1_chem_s4_q04_1}$\xrightarrow[\mathrm{I_2}]{\mathrm{KOH}}$ number of iodoform molecules produced per molecule.\par
(d) Position of ethyl group in correct IUPAC name of following compound.\par
\qfig{p1_chem_s4_q04_2}
Write answer of part (a), (b), (c) \& (d) in the same order and present the four digit number as answer .\par
For example : If all these answer are 9 then fill 9999 .
\qn{5} (a) $\mathrm{CH_3CHO}+\mathrm{CH_3COCH_3}\xrightarrow[5\,^\circ\mathrm{C}]{\text{dil.KOH}}$ is the number of all aldol products. (Excluding stereo)\par
(b) $\mathrm{Ph-C(=O)-CH_3}\xrightarrow[\text{dry ether}]{\text{Mg/Hg}}(A)\xrightarrow[\Delta]{\mathrm{Conc.H_2SO_4}}(B)$\par
Double bond equivalent value of product $B$ will be.\par
(c) How many moles of ethane are released when two moles of $\mathrm{CH_3COONa}$ undergoes Kolbe's electrolytic method.\par
(d) How many monochloro alkenes will give propene on treatment with Mg followed by EtOH. (Including stereo)\par
Write answer of part (a), (b), (c) \& (d) in the same order and present the four digit number as answer .\par
For example : If all these answer are 9 then fill 9999 .
\qn{6} Suggest sequence of 4 reagents in order, from list given below to carry out following conversion strictly in four steps.\par
$\mathrm{Ph-CH=CH-Ph}\to$ Benzene\par
(1) Ammonical $\mathrm{AgNO_3}$ \quad (2) $\mathrm{HIO_4}.\Delta$ \quad (3) (i) NaOH / CaO (ii) $\mathrm{H^+}$ \quad (4) $\mathrm{H^{\oplus}}$, $\mathrm{KMnO_4}$, $\Delta$ \quad (5) Fehling solution \quad (6) KOH, $\mathrm{Br_2}$ \quad (7) $\mathrm{CF_3CO_3H}$ \quad (8) Bayer's reagent
\partsec{3}{MATHEMATICS}
\sect{SECTION-I(i)}
\qn{1} Let $f(x)=x^2+xg'(1)+g''(2)$ and $g(x)=f(1)x^2+xf'(x)+f''(x)$, then the value of $g'(1)+g''(2)$ is -
\begin{qoptions}
\opt{A}{3}
\opt{B}{$-3$}
\opt{C}{0}
\opt{D}{2}
\end{qoptions}
\qn{2} $\displaystyle\int\dfrac{\sin2x+\sin4x-\sin6x}{1+\cos2x+\cos4x+\cos6x}dx$ is equal to (where $C$ is constant of integration)
\begin{qoptions}
\opt{A}{$\dfrac13\ell n|\sec3x|+\dfrac12\ell n|\sec2x|+\ell n|\sec x|+C$}
\opt{B}{$\dfrac13\ell n|\sec3x|-\dfrac12\ell n|\sec2x|-\ell n|\sec x|+C$}
\opt{C}{$\dfrac13\ell n|\sec3x|-\dfrac12\ell n|\sec2x|+\ell n|\sec x|+C$}
\opt{D}{$\dfrac13\ell n|\sec3x|+\dfrac12\ell n|\sec2x|-\ell n|\sec x|+C$}
\end{qoptions}
\qn{3} $f(x)=\tan(\sin^{-1}(\cos(\sin^{-1}(x))))$ $\tan(\cos^{-1}(\sin(\cos^{-1}x)))$ where $x\in(0,1)$ then $\displaystyle\sum_{r=0}^{\infty}\dfrac1{2^r}f\!\left(\dfrac1{2^r}\right)$ is equal to
\begin{qoptions}
\opt{A}{$\dfrac12$}
\opt{B}{$\dfrac52$}
\opt{C}{3}
\opt{D}{2}
\end{qoptions}
\qn{4} Let a real valued differentiable function $f(x)$ be strictly monotonic on $[a,b]$, then
\begin{qoptions}
\opt{A}{{\footnotesize$\displaystyle\int_a^b\left[(f(x))^2-(f(b))^2\right]dx=\int_{f(a)}^{f(b)}2x\left(b-f^{-1}(x)\right)dx$}}
\opt{B}{{\footnotesize$\displaystyle\int_a^b\left[(f(x))^2-(f(a))^2\right]dx=\int_{f(a)}^{f(b)}2x\left(b-f^{-1}(x)\right)dx$}}
\opt{C}{$\displaystyle\int_a^b[f(x)-f(a)]^2dx=\int_{f(a)}^{f(b)}x\left(b-f^{-1}(x)\right)^2dx$}
\opt{D}{None of these}
\end{qoptions}
\sect{SECTION-I(ii)}
\qn{5} Let $f(x)=\displaystyle\int_x^{x+\frac\pi3}|\sin\theta|\,d\theta$; where $x\in[0,\pi]$, which of the following is/are true
\begin{qoptions}
\opt{A}{maximum value of $f(x)$ is 1}
\opt{B}{minimum value of $f(x)$ is $\dfrac32$}
\opt{C}{maximum value of $f(x)$ is $\dfrac12$}
\opt{D}{minimum value of $f(x)$ is $2-\sqrt3$}
\end{qoptions}
\qn{6} Let $f:(a,b)\to R$ is a differentiable function such that $\displaystyle\lim_{x\to a^+}f^2(x)=0$, $\displaystyle\lim_{x\to b^-}f^2(x)=e-1$ and $2f(x)f'(x)-f^2(x)\ge1$ for all $x\in(a,b)$ then value of $(b-a)$ can be
\begin{qoptions}
\opt{A}{0}
\opt{B}{$\dfrac12$}
\opt{C}{1}
\opt{D}{2}
\end{qoptions}
\qn{7} Let $f(x)=\cos^{-1}\!\left(\ell n\!\left(\dfrac{x+1}{x}\right)\right)$, then which of the following statement(s) is/are correct
\begin{qoptions}
\opt{A}{Domain of $f(x)$ contains $(-\infty,-1)\cup(1,\infty)$}
\opt{B}{Range of $f(x)$ contains $\left[0,\dfrac1{e-1}\right]$}
\opt{C}{Domain of $f(x)$ is $\left(-\infty,\dfrac e{1-e}\right]\cup\left[\dfrac1{e-1},\infty\right)$}
\opt{D}{Range of $f(x)$ is $(-\infty,-1]\cup(1,\infty)$}
\end{qoptions}
\sect{SECTION-I(iii)}
\qn{8} Match List-I with List-II and select the correct answer using the code given below the list.\par
{\footnotesize\setlength{\tabcolsep}{2pt}\begin{tabular}{|c|p{0.55\linewidth}|c|p{0.08\linewidth}|}
\hline
\multicolumn{2}{|c|}{\textbf{List-I}} & \multicolumn{2}{c|}{\textbf{List-II}}\\ \hline
(P) & If $x=\alpha$ satisfies $\sin^{-1}x+\cos^{-1}x^2+\dfrac\pi2=0$, then the value of $\dfrac{\sec^{-1}\alpha-\tan^{-1}\alpha}{\cot^{-1}\alpha-\mathrm{cosec}^{-1}\alpha}$ is & (1) & 0\\ \hline
(Q) & If $a=\displaystyle\lim_{x\to\infty}(2.3^x+3.4^x+4.5^x)^{\frac1x}$ and $b=\displaystyle\lim_{x\to(-\infty)}(2.3^x+3.4^x+4.5^x)^{\frac1x}$ then $\left(\dfrac{a+b}{a-b}\right)$ is equal to & (2) & 4\\ \hline
(R) & Number of points of discontinuity of $f(x)=x\sin(\pi(2x-[x]))$ in $[0,2]$ will be (Note $[k]$ denotes greatest integer function less than or equal to $k$) & (3) & 1\\ \hline
(S) & $\displaystyle\lim_{x\to0}\left(\dfrac{3e^{x/3}-x-3}{x^2}\right)$ is equal to & (4) & $\dfrac13$\\ \hline
 & & (5) & $\dfrac16$\\ \hline
\end{tabular}}\par
\begin{qoptions}
\opt{A}{$P\to1;Q\to2;R\to2;S\to4$}
\opt{B}{$P\to3;Q\to1;R\to2;S\to5$}
\opt{C}{$P\to3;Q\to2;R\to1;S\to5$}
\opt{D}{$P\to1;Q\to3;R\to2;S\to4$}
\end{qoptions}
\qn{9} For a continuous function $f(x)$, equation $f(x)=0$ has a 4 distinct solutions then match List-I with List-II and select the correct answer using the code given below the list.\par
{\footnotesize\setlength{\tabcolsep}{2pt}\begin{tabular}{|c|p{0.27\linewidth}|c|p{0.36\linewidth}|}
\hline
\multicolumn{2}{|c|}{\textbf{List-I}} & \multicolumn{2}{c|}{\textbf{List-II}}\\ \hline
(P) & Maximum number of distinct solutions of $f(|x|)=0$ is equal to the value of & (1) & {\scriptsize$\lim\limits_{x\to0^+}\left(\frac{x+1}{2x^2+3x+5}\right)^{1/x}$}\\ \hline
(Q) & Minimum number of distinct solutions of $f(|x|)=0$ is equal to the value of & (2) & {\scriptsize$\lim\limits_{x\to0}\left[\frac{9\sin x}{x}\right]$ (where $[.]$ is greatest integer function)}\\ \hline
(R) & Maximum number of distinct solution of $f(-|x|)=0$ is equal to the value of & (3) & {\scriptsize$\displaystyle\int_2^{18}\frac{f(x)}{f(20-x)+f(x)}dx$}\\ \hline
(S) & Minimum number of distinct solution of $f(-|x|)=0$ is equal to the value of & (4) & {\small$\displaystyle\int_{-2}^{2}\dfrac{x^3-x}{\sqrt{4-x^2}}dx$}\\ \hline
 & & (5) & {\small$24\displaystyle\int_{1/3}^{2/3}\dfrac{\ell nx}{\ell n(x-x^2)}dx$}\\ \hline
\end{tabular}}\par
\begin{qoptions}
\opt{A}{$P\to3,4;Q\to1,2,4;R\to2,3;S\to1,2,4$}
\opt{B}{$P\to2,5;Q\to1,4,5;R\to2,3;S\to1,2,4$}
\opt{C}{$P\to2,3;Q\to1,4;R\to2,3;S\to1,4$}
\opt{D}{$P\to1,2,3;Q\to1,3;R\to1,2,3;S\to1,2,4$}
\end{qoptions}
\qn{10} Match List-I with List-II and select the correct answer using the code given below the list.\par
{\footnotesize\setlength{\tabcolsep}{2pt}\begin{tabular}{|c|p{0.5\linewidth}|c|p{0.08\linewidth}|}
\hline
\multicolumn{2}{|c|}{\textbf{List-I}} & \multicolumn{2}{c|}{\textbf{List-II}}\\ \hline
(P) & Value of $\dfrac8{\pi^2}\displaystyle\int_0^\pi\sin^{-1}(\sin x)\,dx$ is equal to & (1) & 0\\ \hline
(Q) & Value of $\displaystyle\int_0^\pi\cos mx\ \sin nx\,dx$ is equal to (Here $m$ and $n\in I$ \& $(m-n)$ is a non-zero even integer) & (2) & 1\\ \hline
(R) & Value of $\displaystyle\lim_{x\to0^+}\sin\cos^{-1}\cot^{-1}\!\left(\dfrac1x\right)$ is equal to & (3) & 2\\ \hline
(S) & $\displaystyle\lim_{x\to2^+}\left(\dfrac{2\{x\}-4}{[x]-3}\right)$ is equal to (Here $[.]$ \& $\{.\}$ denote greatest integer \& fractional part function) & (4) & 4\\ \hline
 & & (5) & 3\\ \hline
\end{tabular}}\par
\begin{qoptions}
\opt{A}{$P\to3;Q\to2;R\to1;S\to4$}
\opt{B}{$P\to2;Q\to3;R\to5;S\to1$}
\opt{C}{$P\to3;Q\to1;R\to2;S\to4$}
\opt{D}{$P\to5;Q\to2;R\to4;S\to1$}
\end{qoptions}
\qn{11} Match the statement in List-I with given value in List-II\par
{\footnotesize\setlength{\tabcolsep}{2pt}\begin{tabular}{|c|p{0.62\linewidth}|c|p{0.06\linewidth}|}
\hline
\multicolumn{2}{|c|}{\textbf{List-I}} & \multicolumn{2}{c|}{\textbf{List-II}}\\ \hline
(P) & If area bounded by $x=(\theta-\sin\theta)$; $y=(1-\cos\theta)$; $0\le\theta\le2\pi$ is $\lambda\pi$; then the value of $\lambda$ can be equal to & (1) & 6\\ \hline
(Q) & If the least area of triangle formed by tangent to the circle $x^2+y^2=1$ and $x=0$, $y=0$ is $A$, then $A$ is coprime with & (2) & 3\\ \hline
(R) & Area of triangle formed by any tangent to the curve $y=\dfrac{2x^2-x-1}{2x^2+5x+2}\left(x\ne-2,-\dfrac12\right)$ and pair of lines $xy-x+2y=2$ is & (3) & 80\\ \hline
(S) & If the maximum area bounded by the curves $x^2=4ay$, $y=ax$, $y=\dfrac xa$, $1\le a\le2$ is $\alpha$, then value of $(\alpha+a)$ is & (4) & 84\\ \hline
 & & (5) & 86\\ \hline
\end{tabular}}\par
\begin{qoptions}
\opt{A}{$P\to2;Q\to1,2;R\to1,4,5;S\to1,2$}
\opt{B}{$P\to2;Q\to1,2,3,4,5;R\to1;S\to4$}
\opt{C}{$P\to2;Q\to1,2,3,4,5;R\to2,5;S\to1,4$}
\opt{D}{$P\to3;Q\to2;R\to4;S\to4$}
\end{qoptions}
\sect{SECTION-II}
\qn{1} Given $\displaystyle\sum_{n=1}^{\infty}\dfrac1{n^2}=\dfrac{\pi^2}6$ and $\displaystyle\int_0^1\dfrac{\log_ex}{1-x^2}dx=-\dfrac{\pi^2}{\lambda}$, then the value of $\lambda$ equals
\qn{2} The value of $\displaystyle\lim_{n\to\infty}\dfrac1n\sum_{k=1}^{n}\left[e^{(1+k/n)^2}-\dfrac{3e^{\left(1+\frac{3k}n\right)}}{2\sqrt{1+\frac{3k}n}}\right]$ equals
\qn{3} Let $f(x)=\dfrac{x}{1+(\ell nx)(\ell nx)(\ell nx)\ldots\infty}\ \forall x\in[1,3]$, is non-differentiable at $x=k$. Then the value of $[k^2]$, is ($[.]$ denotes the greatest integer function)
\qn{4} Let $C_1$, $C_2$ and $C_3$ be the three circles lying in first quadrant each touching both axes. $C_1$ touches $C_2$ and $C_2$ touches $C_3$. If area of smallest circle is 1 and area of largest circle is $\left(a+2\sqrt b\right)^4$, (where $a$ and $b$ are coprime) then value of $a+b$ is
\qn{5} A circle of radius $2\sqrt{10}$ cm is divided by chord of length 12cm into two parts. If maximum area of rectangle inscribed into smaller circular segment is $\lambda$, then $\dfrac\lambda{\sqrt{15}}$ is
\qn{6} The value of\par{\small$\displaystyle\lim_{x\to\frac\pi4}\left(\dfrac{4\sqrt2-(\cos x+\sin x)^5-10\cos x+10\cos x\sin2x}{1-\sin2x}\right)$} is
\papersec{2}
\partsec{1}{PHYSICS}
\sect{SECTION-I(i)}
\qn{1} There is a shallow cylindrical container of radius $R$ and height $\dfrac R2$ as in the figure filled completely with a coloured liquid. The container is rotated such that $\dfrac14$th area of the bottom is almost dry. The force exerted by liquid on the bottom of the container in this position is :-\par
(Neglect atmospheric pressure, $\rho$ = density of liquid)
\qfig{p2_phy_s1_q01_1}
\begin{qoptions}
\opt{A}{$\rho g\dfrac{R^2}{2}$}
\opt{B}{$\dfrac1{12}\rho gR^2$}
\opt{C}{$\dfrac{23}{128}\rho gR^2$}
\opt{D}{$\dfrac{45}{256}\rho gR^2$}
\end{qoptions}
\qn{2} An alternating current flows as in fig.(I) in circuit of fig.(II). Find r.m.s. of voltage across $A$ and $B$.
\qfig{p2_phy_s1_q02_1}
\qfig{p2_phy_s1_q02_2}
\begin{qoptions}
\opt{A}{$i_0R$}
\opt{B}{$i_0\sqrt{R^2+\dfrac{L^2}{T^2}}$}
\opt{C}{$i_0\sqrt{R^2+\dfrac{4\pi^2L^2}{T^2}}$}
\opt{D}{$i_0\sqrt{\dfrac{R^2}3+\dfrac{4L^2}{T^2}}$}
\end{qoptions}
\qn{3} A 1 kg mass stretches a spring by 49 cm. The system is immersed in a medium that impacts a damping force equal to four times the instantaneous velocity of the mass. Particle is released from a point 24 cm above the equilibrium. If upward direction is taken positive then equation of motion is -
\begin{qoptions}
\opt{A}{$e^{-4t}\left[0.24\cos4t+0.24\sin4t\right]$}
\opt{B}{$e^{-2t}\left[0.24\cos4t+0.24\sin4t\right]$}
\opt{C}{$e^{-2t}\left[0.24\cos4t+0.12\sin4t\right]$}
\opt{D}{$e^{-4t}\left[0.24\cos4t-0.24\sin4t\right]$}
\end{qoptions}
\qn{4} A block is connected to a spring. The spring is extended when block is released. The block experiences fluid friction and executes damped S.H.M. Energy of S.H.M. depends on time. Which of the following graph most closely represent this variation ?
\qfig{p2_phy_s1_q04_1}
\begin{qoptions}
\opt{A}{\qfigi{p2_phy_s1_q04_A}}
\opt{B}{\qfigi{p2_phy_s1_q04_B}}
\opt{C}{\qfigi{p2_phy_s1_q04_C}}
\opt{D}{\qfigi{p2_phy_s1_q04_D}}
\end{qoptions}
\sect{SECTION-I(ii)}
\qn{5} In the figure there is a conducting shell of radius $R$ and thickness $t$ ($\ll R$) in which a current $I$ enters and leaves at diametrical opposite points. ``$V$'' is the potential difference between $A$ and $B$. Choose the correct option.
\qfig{p2_phy_s2_q05_1}
\begin{qoptions}
\opt{A}{$V$ is directly proportional to $I$.}
\opt{B}{$V$ is inversely proportional to $R^3$.}
\opt{C}{$V$ is directly proportional to $t$.}
\opt{D}{$V$ is inversely proportional to $t$.}
\end{qoptions}
\qn{6} In the figure shown both region I and II contain uniform magnetic field $B_0$ in opposite direction and charge particle of charge ``$q$'' and mass $m$ is thrown radially from common centre then :-
\qfig{p2_phy_s2_q06_1}
\begin{qoptions}
\opt{A}{If $v<\dfrac{qB_0R}{2m}$, charge will remain in region I.}
\opt{B}{If $v>\dfrac{qB_0R}{2m}$, charge will enter in region-III.}
\opt{C}{If $v>\dfrac{qB_0R}{2m}$, charge will enter in region-III radially, no matters how large is $v$.}
\opt{D}{Charge will enter region-III with speed more than $v$.}
\end{qoptions}
\qn{7} Consider a cylinder of radius $R$ and height $H$ moving in air such that the speed of center of cylinder is $v_{cm}$ and the angular velocity of rotation about the axis of cylinder is $\omega$.
\qfig{p2_phy_s2_q07_1}
The layer of air at the surface of cylinder is swept away due to rotation of the cylinder the velocity of air at the surface can be considered the resultant of $\vec v_A$ and $\vec v_P$, where $\vec v_A$ and $\vec v_P$ are velocities of air and point $P$ on the cylinder as seen from a frame moving with a velocity $v_{cm}$ towards positive $x$-direction. Assuming laminar flow and pressure at forward point $A$ to be $P_0$.
\begin{qoptions}
\opt{A}{The pressure at point $P$ on the circumference of the cylinder is $p=p_0-v_B\omega R\rho_{\text{air}}\sin\theta$}
\opt{B}{The magnitude of force acting on the cylinder due to air flow is $\rho_{\text{air}}\omega v_{cm}\pi R^2L$.}
\opt{C}{The direction of force acting on the cylinder due to air flow is $+y$ direction.}
\opt{D}{None of the above.}
\end{qoptions}
\sect{SECTION-II(i)}
\paraq{1}{2}
A pulse is started at time $t=0$ along the $+x$ direction on a long uniform string of linear mass density 10 gm/m at tension 100 N. The shape of the pulse at $t=0$ is given by :\par
$y=\begin{cases}\dfrac x4+1 & -4<x\le0\\ \sqrt3(1-x) & 0<x<1\\ 0 & \text{Otherwise}\end{cases}$\par
Here $x$ and $y$ are in centimeter.
\qn{1} At the same instant speed of particle at $x=-2$cm in m/s is :-
\qn{2} At this instant magnitude of rate of transmission of energy at $x=\dfrac12$ cm is $x\times10^4$ W. Find $x$.
\paraq{3}{4}
The electric fields of two plane electromagnetic plane waves in vacuum are given by\par
$\vec E_1=E_{0}\hat j\cos(\omega t-kx)$ and $\vec E_2=E_0\hat k\cos(\omega t-ky)$\par
At $t=0$, a particle of charge $q$ is at origin moving with a velocity $\vec v=0.75c\,\hat j$. $c$ is the speed of light.
\qn{3} Find the ratio of net force experienced and electric force on charge due to first wave only at the same instant.
\qn{4} If net force due to both the waves at the same instant is $\vec F=qE_0\left[a\hat i+b\hat j-\dfrac c2\hat k\right]$. Find $a+b+c$.
\sect{SECTION-II(ii)}
\qn{5} Circular region has unsteady magnetic field $\vec B=\alpha t\hat k$. $ACB$ is a conducting circular arc of same radius. Potential difference between $A$ and $B$ is $\alpha R^a\left(1-\dfrac{\sqrt b}{\sqrt c}\right)$. Find $a+b+c$. (Use : $\pi=3$)
\qfig{p2_phy_s4_q05_1}
\qn{6} Figure shows a uniform solid cylinder released from the position shown. If cylinder is just able to come out of water completely. Find the ratio of density of water to density of cylinder.
\qfig{p2_phy_s4_q06_1}
\qn{7} A sound source is moving as shown in figure for which circles drawn are wavefront with a phase difference of $2\pi$. Find distance $AB$. Given velocity of source $=12.5$ m/s. Velocity of sound $=330$ m/s, Frequency emitted $=160$ Hz.
\qfig{p2_phy_s4_q07_1}
\qn{8} If current is marked $2\,\Omega$ resistance is 2A. Find unknown resistance $R$ in $\Omega$.
\qfig{p2_phy_s4_q08_1}
\qn{9} Figure shows a conducting ball of capacitance $2\mu$F colliding with two large parallel plates elastically. The equivalent current in mA is :- (Assume ball moves with constant speed as shown in figure)
\qfig{p2_phy_s4_q09_1}
\qn{10} A transmission line consists of two concentric tubes of thin copper sheet metal with radii $R_1=R$ and $R_2=2R$ (see figure). The current flows to the left along the central tubes and back along the other, completing a circuit; the current is uniformly distributed over the surface of each tube. The self-inductance per meter length of this transmission line is $\left(\dfrac1{10}\right)\ell n(X)$ micro henry ($\mu$H). Find the value of $X$.
\qfig{p2_phy_s4_q10_1}
\partsec{2}{CHEMISTRY}
\sect{SECTION-I(i)}
\qn{1} Assuming that during the formation of $\mathrm{X_2^+}$ from $\mathrm{X_2}$ the bond length increase while during the formation of $\mathrm{M_2^{-2}}$ from $\mathrm{M_2}$ bond length decrease, where ($6<$ atomic number of $x<10$) and ($5<$ atomic number of $M<10$)\par
From given information which of the following statement is incorrect?
\begin{qoptions}
\opt{A}{During the formation of $\mathrm{X_2^+}$ from $\mathrm{X_2}$, $e^-$ is removed from gerade orbital.}
\opt{B}{During formation of $\mathrm{M_2^{-2}}$ from $\mathrm{M_2}$, $e^-$ is added to ungerade orbital.}
\opt{C}{Total nodal plane in that molecular orbital from which $e^-$ is removed from $\mathrm{X_2}$ to form $\mathrm{X_2^+}$ is 2}
\opt{D}{Total nodal plane in that molecular orbital in which last $e^-$ of $\mathrm{M_2}$ is available is 2.}
\end{qoptions}
\qn{2} Which of the following statement is correct?
\begin{qoptions}
\opt{A}{Except alkali metal carbonate and ammonium carbonate , all normal carbonates are insoluble in water}
\opt{B}{Lime water and $\mathrm{Cr_2O_7^{-2}}$ solution both are used to differentiate $\mathrm{CO_2}$ and $\mathrm{SO_2}$ gas.}
\opt{C}{Iodine Azide test is specific test of Sulphide ion only.}
\opt{D}{Acetate ion, gives white ppt. with $\mathrm{FeCl_3}$.}
\end{qoptions}
\qn{3} The enthalpy of atomization of $\mathrm{P_4O_6(s)}$ is 3500 kJ/mol. If P--O bond enthalpy is 250 kJ/mol then enthalpy of sublimation of $\mathrm{P_4O_6(s)}$ is :
\begin{qoptions}
\opt{A}{500 kJ/mol}
\opt{B}{600 kJ/mol}
\opt{C}{300 kJ/mol}
\opt{D}{200 kJ/mol}
\end{qoptions}
\qn{4}
\qfig{p2_chem_s1_q04_1}
Reagents $R_1$, $R_2$ and $R_3$ should be respectively.\par
\begin{qoptions}
\opt{A}{$\mathrm{NH_2{-}NH_2,\ OH^-}$;\ \ $\mathrm{Zn{-}Hg/conc.HCl}$;\ \ Se, $\Delta$}
\opt{B}{$\mathrm{Zn{-}Hg/conc.\ HCl}$;\ \ $\mathrm{NH_2{-}NH_2,\ OH^-}$;\ \ Se, $\Delta$}
\opt{C}{$\mathrm{Zn{-}Hg/conc.\ HCl}$;\ \ $\mathrm{NH_2{-}NH_2,\ OH^-}$;\ \ Cu, 300$^\circ$C}
\opt{D}{$\mathrm{NH_2{-}NH_2,\ OH^-}$;\ \ $\mathrm{Zn{-}Hg/conc.\ HCl}$;\ \ Cu, 300$^\circ$C}
\end{qoptions}
\sect{SECTION-I(ii)}
\qn{5} Which of the following expressions are correct for an ideal gas :
\begin{qoptions}
\opt{A}{$\left(\dfrac{\partial U}{\partial V}\right)_T=0$}
\opt{B}{$\left(\dfrac{\partial U}{\partial P}\right)_T=0$}
\opt{C}{$\left(\dfrac{\partial S}{\partial V}\right)_T=0$}
\opt{D}{$\left(\dfrac{\partial G}{\partial P}\right)_T=0$}
\end{qoptions}
\qfullbegin
\qn{6}
\qfig{p2_chem_s2_q06_1}
\qfullsplit
The correct statement about above reaction sequence.
\begin{qoptions}
\opt{A}{`P' and `Q' both give $+$ve Iodoform test.}
\opt{B}{$\mathrm{LiAlH_4}$ reduces carbonyl compound to an alcohol.}
\opt{C}{`S' is an unsaturated alcohol (with DU $=6$)}
\opt{D}{`R' is \qfigi{p2_chem_s2_q06_2}}
\end{qoptions}
\qfullend
\qn{7} Which of the following reaction are correctly matched?
\begin{qoptions}
\opt{A}{\qfigi{p2_chem_s2_q07_A}}
\opt{B}{\qfigi{p2_chem_s2_q07_B}}
\opt{C}{\qfigi{p2_chem_s2_q07_C}}
\opt{D}{$\mathrm{CH_3{-}O{-}CH_2{-}CH_3}\xrightarrow{\mathrm{HBr}}S_N2$}
\end{qoptions}
\sect{SECTION-II(i)}
\paraq{1}{2}
The crystal field splitting, $\Delta_o$ depends upon the field produced by the ligand and charge on the metal ion. Some ligands are able to produce strong fields in which case, the splitting will be large whereas others produce weak field and consequently result in small splitting of d-orbitals.
\qfig{p2_chem_s3_para1_2_1}
\qn{1} Maximum number of atoms which can lie in same plane in complex $B$ is :-
\qn{2} Sum of total number of stereoisomers in complex ($C$) is $M$ and number of unpaired $e^-$ in central atom of complex ($A$) is $N$ then find $(M+N)^2$
\paraq{3}{4}
An analyst starts a first order chemical reaction at 8.00 A.M. in the morning at the laboratory temperature of $27^\circ$C. At 1.00 P.M., he discovered that only 10\% of the reaction was complete by that time. To speed-up the reaction he increased the temperature to $127^\circ$C. At 4.00 P.M. he found that only 50\% of the reaction was complete from begining. Any how he did not want to stay in laboratory beyond 5.00 P.M. but he could not leave the laboratory until the reaction was 90\% complete. Fortunately he found a suitable catalyst, adding which at 4.00 P.M. at $127^\circ$C, he could meet the target of 5.00 P.M. Answer the following questions based on the above observation.\par
(Use : $\ln\dfrac{10}9=0.1$, $\ln\dfrac95=0.6$ , $\ln10=2.3$ , $\ln5=1.6$, $\ln8=2$)
\qn{3} What was the activation energy of the original pathways (in kcal/mol) ?
\qn{4} What was the activation energy of the catalyzed pathway (in kcal/mol) ?
\sect{SECTION-II(ii)}
\qn{5} How many following statements is/are correct?\par
(i) In Mond's process a complex is formed before the reduction of metal involves d-orbital having zero nodal plane in hybridisation of central atom\par
(ii) In Cyanide process of silver glance disproportionation redox reaction take place \& $n$-factor of sulphur atom is $\dfrac85$.\par
(iii) Tempering is process in which steel becomes quite hard and brittle.\par
(iv) In electrorefining of Aluminium molten mixture of Cryolite and Fluspar is used as electrolyte.
\qn{6} In the following ores find sum of $(x+y+z)$ where\par
$x$ = number of ores in which metals exist in different oxidation state.\par
$y$ = number of ores in which two different metals are available.\par
$z$ = number of ores in which atleast one metal ion is diamagnetic in nature.\par
Magnetite, Chromite, Limonite, Malachite, I$\ell$menite, Chalcopyrite, Dolomite, Azurite, Carnalite, Cryolite
\qn{7} If each hydrogen atom in the ground state, 1.0 mol of H atoms are excited by absorbing photons of energy 8.4 eV, 12.09 eV, and 15.0 eV of energy, then the number of different types of spectral lines emitted is equal to.
\qn{8} Find freezing point temperature (in Kelvin)of liquid if vapour pressure of solid \& liquid are given by following expression.\par
$\ell n p_s=10-\dfrac{3000(\mathrm K)}{T}$ \quad where $P_s$ = vapour pressure of solid\par
$\ell n p_\ell=5-\dfrac{2000(\mathrm K)}{T}$ \quad $P_\ell$ = vapour pressure of liquid\par
\textbf{Note :} All terms are in SI unit
\qn{9} Consider the following molecule.
\qfig{p2_chem_s4_q09_1}
$P_1$ : moles of $\mathrm{H_2}$ evolved after reaction with Na.\par
$P_2$ : number of site where Lucas reagent will give turbidity.\par
$P_3$ : number of site where neutral $\mathrm{FeCl_3}$ will react\par
$P_4$ : moles of $\mathrm{CO_2}$ evolved when 1 mole of compound reacts 1 mole of $\mathrm{NaHCO_3}$.\par
$P_5$ : number of site where Tollen's reagent will react (either Ag mirror or white precipitate).\par
$P_6$ : number of site where 2,4-DNP will react.\par
Write the sum of $P_1+P_2+P_3+P_4+P_5+P_6$.
\qn{10}
\qfig{p2_chem_s4_q09_2}
(a) The number of carbon atoms present in the principal carbon chain of product $P_3$ obtained when $P_1$ is heated with sodalime.\par
(b) How many moles of gas $P_2$ released when one mole of $P_1$ is treated with sufficient amount of $\mathrm{NaHCO_3}$ ?\par
(c) How many moles of $\mathrm{CH_3COCl}$ reacts with 1 mole of $P_4$ ?\par
(d) What is D.B.E. (Double bond equivalent) for $P_1$ ?\par
Write answer of part (a), (b), (c) \& (d) in the same order and present the four digit number as answer .\par
For example : If all these answer are 9 then fill 9999 .
\partsec{3}{MATHEMATICS}
\sect{SECTION-I(i)}
\qn{1} The absolute value of the expression $y^3\dfrac{d^2y}{dx^2}$ for the ellipse $3x^2+4y^2=12$ is $\dfrac PQ$ (where $P$ and $Q$ are coprime) then the value of $(P-Q)$ is
\begin{qoptions}
\opt{A}{2}
\opt{B}{3}
\opt{C}{4}
\opt{D}{5}
\end{qoptions}
\qn{2} Let $f(x)=\displaystyle\lim_{h\to0}\dfrac{(\sin(x+h))^{\ell n(x+h)}-(\sin x)^{\ell nx}}{h}$ then $f\!\left(\dfrac\pi2\right)$ is
\begin{qoptions}
\opt{A}{0}
\opt{B}{1}
\opt{C}{$\ell n\dfrac\pi2$}
\opt{D}{Does not exist}
\end{qoptions}
\qn{3} A circle touches two perpendicular lines $OA$ and $OB$ at points $M$ and $N$ respectively and passes through the vertex $C$ of rectangle $OACB$ (in this order). If the perpendicular distance of $C$ from the line $MN$ is of 10 units length, then the area of the rectangle $OACB$ is
\begin{qoptions}
\opt{A}{50}
\opt{B}{100}
\opt{C}{225}
\opt{D}{25}
\end{qoptions}
\qn{4} If $C_1\equiv y=\dfrac1{1+x^2}$ and $C_2\equiv y=\dfrac{x^2}2$ be two curve lying in $XY$ plane. Then
\begin{qoptions}
\opt{A}{area bounded by curve $y=\dfrac1{1+x^2}$ and $y=0$ is $2\pi$}
\opt{B}{area bounded by $C_1$ and $C_2$ is $\dfrac\pi2-\dfrac13$}
\opt{C}{area bounded by $C_1$ and $C_2$ is $1-\dfrac\pi2$}
\opt{D}{area bounded by curve $y=\dfrac1{1+x^2}$ and $x$-axis is $\dfrac\pi2$}
\end{qoptions}
\sect{SECTION-I(ii)}
\qn{5} Let a differentiable function $f:R\to R$ be such that for all $x$ and $y$ in $R$ we have $2|f(x)-f(y)|\le|x-y|$ and $f'(x)\geqslant\dfrac12$. Then number of points of intersection the graph of $y=f(x)$ with
\begin{qoptions}
\opt{A}{the line $y=x$ is one}
\opt{B}{the curve $y=-x^3$ is one}
\opt{C}{the curve $2y=|x|$ can be zero}
\opt{D}{the curve $y^2=-x$ may be more than one}
\end{qoptions}
\qn{6} If $p=\displaystyle\lim_{n\to\infty}n^{-n^2}\Big\{(n+2^0)(n+2^{-1})$ $(n+2^{-2})\ldots(n+2^{-n+1})\Big\}^n$ then which of the following is true ?
\begin{qoptions}
\opt{A}{$p$ is irrational}
\opt{B}{integer nearest to $p$ is 7}
\opt{C}{$p>\displaystyle\lim_{x\to\infty}\left(\dfrac{1+x}x\right)^{\ell nx-x}$}
\opt{D}{None of these}
\end{qoptions}
\qn{7} Let $g:R\to R$ be a function defined by $g(x)=x^3+x+4$ and $f$ is a function such that $f(g^{-1}(x))=x^3+1$, then
\begin{qoptions}
\opt{A}{$\left.\dfrac d{dx}(f(x))\right|_{x=0}=48$}
\opt{B}{$\left.\dfrac d{dx}\left(f^{-1}(x)\right)\right|_{x=9}=\dfrac1{48}$}
\opt{C}{$\displaystyle\int_0^1\dfrac{\sqrt{g(x)}}{f(x)}\left(3x^2+1\right)dx=\dfrac23\tan^{-1}\left(\dfrac{6\sqrt6-8}{1+48\sqrt6}\right)$}
\opt{D}{$\displaystyle\int_0^1\dfrac{\sqrt{g(x)}}{f(x)}\left(3x^2+1\right)dx=\dfrac23\tan^{-1}\left(\dfrac2{25}\right)$}
\end{qoptions}
\sect{SECTION-II(i)}
\paraq{1}{2}
Suppose $f(x)$ and $g(x)$ are differentiable function such that $xg(f(x))\cdot f'(g(x))\cdot g'(x)=f(g(x))\cdot g'(f(x))\cdot f'(x)$ $\forall\ x\in R$ and $f(x)$ \& $g(x)$ are positive for all $x\in R$.\par
Also $\displaystyle\int_0^xf(g(t))dt=\dfrac12\left(1-e^{-2x}\right)\forall x\in R$, $g(f(0))=1$ and $h(x)=\dfrac{g(f(x))}{f(g(x))}\forall x\in R$.
\qn{1} If the graph of $y=h(x)$ is symmetric with respect to the line $x=\lambda$, then $\lambda$ equals
\qn{2} If the largest possible value of $h(x)\ \forall\ x\in R$ is $t$, then $[t]$ equals (where $[\cdot]$ denotes greatest integer function)
\paraq{3}{4}
Let $f(x)$ be defined in the interval $(-\infty,\infty)$ such that\par
$f(x)=\begin{cases}\dfrac{a(1-x\sin x)+b\cos x+5}{x^2} & x<0\\[2mm] 3 & x=0\\[1mm] \left(1+\left(\dfrac{cx+dx^3}{x^2}\right)\right)^{\frac1x} & x>0\end{cases}$ and\par
$g(x)$ be defined in the interval $(c,-b)$ such that\par
$g(x)=\begin{cases}\dfrac{x^2+a}{x^2+1} & c<x\leqslant2+c\\[2mm] \dfrac14\left(x^3-x^2+c\right) & 2+c<x\leqslant-b-1\\[2mm] \dfrac94\left(|x+b|+|x+2a|\right) & -b-1<x<-b\end{cases}$\par
On the basis of above information, answer the following questions :
\qn{3} If $f(x)$ is continuous at $x=0$, then $\dfrac{\left|a+b+c+e^d\right|}4$ equals
\qn{4} The number of points of non-differentiability of $g(x)$ in $x\in(c,-b)$
\sect{SECTION-II(ii)}
\qn{5} Let $\lambda=\displaystyle\int_1^9\left(\dfrac12+\sqrt{\dfrac14+\log_3x}\right)dx+\int_1^2\dfrac{3^{x^2-[x]}}{3^{\{x\}}}dx$ then $\lambda$ equals\par
(where [.] denotes greatest integer function and \{.\} denotes fractional part function)
\qn{6} The parabola $y^2=x-y$ divides the triangle formed by the tangent to it at the origin and the straight line $x+y=0$ and $y=-1$ into two parts whose areas are $A_1$ and $A_2$ (Given $A_1<A_2$). If $\dfrac{A_1}{A_2}=\dfrac pq$ where $p,q\in N$ and $\dfrac pq$ is in simplest form then the value of $(p+q)$ is :
\qn{7} Let $L_1$, $L_2$ and $L_3$ be the length of tangents drawn from a point $P$ to the circles $x^2+y^2=4$, $x^2+y^2-4x=0$ and $x^2+y^2-4y=0$ respectively. If $L_1^4=L_2^2L_3^2+16$, then locus of $P$ are the curves $C_1$(a straight line) and $C_2$(a circle). Let the circumcentre of the triangle formed by $C_1$ and other two lines which are inclined at an angle of $45^\circ$ with $C_1$ and is tangent to $C_2$ is $(\alpha,\beta)$, then $(\alpha+2\beta)$ equals
\qn{8} Find the values of $x$ satisfying the equation $2\cos^{-1}(1-x)-3\cos^{-1}x=\pi$
\qn{9} Let $f(x)$ be a differentiable function on $[0,8]$ such that $f(1)=5$, $f(2)=\dfrac12$, $f(3)=6$, $f(4)=-3$, $f(5)=4$, $f(6)=\dfrac14$ and $f(7)=\dfrac{-1}5$. If the minimum number of roots of the equation $f'(x)-f'(x)(f(x))^2=0$ is $\lambda$, then $\lambda$ equals
\qn{10} Let $P(x)$ be a polynomial of degree 4 having a local maximum at $x=2$ and $\displaystyle\lim_{x\to0}\left(3-\dfrac{P(x)}x\right)=27$. If $P(1)=-9$ and $P''(x)$ has a local minimum at $x=2$. If global maximum value of $y=P'(x)$ on the set $A=\{x:x^2+12\le7x\}$ is $4M$ then value of $M$ is :
\end{multicols}
\end{document}
